\section{Contributions}

\section{Limitations}
    The work presented in this thesis has several limitations that must be acknowledged honestly.
    % If no user tests
    \textbf{No empirical user study.} The pedagogical effectiveness of the system has not been evaluated in a controlled experiment.
    \textbf{Limited content library.} The current content library is limited in scope and diversity.

    \textbf{Steep onboarding barriers.} The Glosy same of LingQ, incorporates real-world texts, its interface poses steep onboarding barriers for novice learners \cite{karasimos2022battle}.

    \textbf{Lack of social scaffolding.} The system lacks mechanisms for social scaffolding through peer or native-speaker interaction.

    % Write this claude: some language are formal/informal spoken/written like farsi. English and Greek is not like that. Also some languages have suffixes and prefixes that form grammer like turkish and farsi. for example مادَر is mother and my mother in colloquial is مادَرَم. In English, we don't have that. that's why Tokenizers (hazm) fails to tokenize all words in farsi that are included in out database words table. In the future give option to manually tokenize words while creating a reel. 

\section{Future Work}
    The following directions are identified as high-value:
    \begin{sloppypar}
    \begin{itemize}
        \item \textbf{Automated content annotation pipeline:} A pipeline that ingests a raw video URL, transcribes speech using a large language model, aligns the transcript to video timestamps, tokenises and POS-tags the output, maps tokens to the word database, and produces a fully annotated reel record—reducing the per-reel curation cost from hours to minutes.
        \item \textbf{Empirical evaluation:} A randomised controlled study measuring vocabulary gain, retention at one week and one month, and learner engagement metrics, comparing Glosy against a no-game control condition and against Duolingo.
        \item \textbf{Multi-language expansion:} The schema and service layer are language-agnostic (language is a foreign key throughout); adding a new language requires only content ingestion, not code changes.
        \item \textbf{Collaborative learning features:} Shared vocabulary challenges, leaderboards scoped to friend groups, and comment-driven vocabulary discussion on individual reels, supporting the social dimension of Vygotsky's theory.
        \item \textbf{Vocabulary synchronisation performance:} Replacing the serial per-word UPDATE loop in \texttt{syncController.js} with a bulk \texttt{INSERT ... ON CONFLICT DO UPDATE} statement to support large offline sessions without latency degradation.
        % use Cache like Redis
        \item \textbf{Cold-start mitigation via onboarding preferences:} Wiring the topic preferences and age already collected by onboarding's \texttt{PersonalizationSlide.tsx} (Section~\ref{subsec:recommender_cold_start}) into Stage~3 as a cold-start substitute for the interaction-based profile vector, replacing its current fixed, hand-picked eight-topic taxonomy with automated topic generation derived from reel content once the library is large enough to support it.
        \item \textbf{Age-conditioned recommendation:} Using the \texttt{age} column, already collected at onboarding but currently unread by the recommendation engine, as a further ranking or filtering signal once a concrete policy for it is designed.
    \end{itemize}
    \end{sloppypar}

    
\newpage